%% kapa: exam=1 round=20 year=2009 period=2 form=A subject=appraisal_law subject_ko=부동산관계법규 seq=30 number=70
%% source: Q-Net 감정평가사 기출문제 (cst00309) artlSeq=5024178 file=extracted/1008/5024178/2009년 제20회 감정평가사 1차시험 기출문제/2교시.hwp
%% conversion: kordoc 4.24.51 -> tools/gen1.py (auto); review=proofread-sonnet; flags=none
\begin{question}{70}
건축법령 및 판례에 의할 때 일조권에 관한 설명으로 옳지 \uline{않은} 것은?
\begin{choices}
  \item 동지를 기준으로 오전 9시에서 오후 3시까지 사이에 일조시간을 2시간 이상 계속하여 확보하면 수인한도를 넘지 않는 것으로 본다.
  \item 동지를 기준으로 오전 8시부터 오후 4시까지 사이의 8시간 중 일조시간을 통틀어 4시간 이상 확보하면 수인한도를 넘지 않는 것으로 본다.
  \item 주거의 일조는 쾌적하고 건강한 생활에 필요한 생활이익으로서 가치가 있다고 인정되면 법적 보호의 대상이 된다.
  \item 초등학교 학생과 같이 건물을 일시적으로 이용하는 자도 생활이익으로서의 일조권을 법적으로 보호받을 수 있는 지위에 있다.
  \item 일조권의 침해를 받았더라도 건축허가에 기하여 이미 건축공사를 완료하였다면 그 건축허가처분의 취소를 구할 소의 이익이 없다.
\end{choices}
\end{question}
\answer{4}
